\section{Discussion and Open Problems}
\label{sec:discussion-body}

Three questions suggest themselves. First, whether the bound of
Theorem~\ref{thm:gap} is order-optimal for all nilpotent classes; our
Proposition~\ref{prop:nilpotent} settles only class $\le 2$. Second,
whether the \emph{bipartite} double cover $\tilde G = G \times \mathbb{Z}/2$
preserves the gap up to a universal constant -- the parity obstruction makes
this delicate. Third, whether sharp constants can be derived from
Kesten-type equations \eqref{eq:rho} on the group algebra.

We expect that a Fourier-analytic approach on the dual $\hat G$, combined
with the transfer principle of \cite{Lubotzky}, yields an upper bound of
the form
\begin{equation}
\label{eq:conj}
1 - \lambda_1 \;\le\; \frac{C}{d} \left( 1 - \frac{\log |G|}{\log 2} \right)^{-1}.
\end{equation}

\begin{quote}
``The spectral radius is the shadow of the group's geometry.'' --- folklore
\end{quote}

Figure~\ref{fig:cone} illustrates the lattice approximation used in the proof
of Proposition~\ref{prop:nilpotent}, and Figure~\ref{fig:spectrum} reports the
numerics mentioned in the introduction. Both figures were produced with
TikZ and matplotlib respectively.
